\chapter{Le signal d'alarme}
% version complète: chapters/14_pain.tex

\pencilfig{14}{\pencilart{14}}{La douleur est une alarme incendie, pas un thermomètre.}

Tous les autres sens disent comment est le monde. La douleur, non. Elle signale que quelque chose menace le corps, et son volume est fixé en grande partie par la machine elle-même, et non par le stimulus. La douleur est une alarme incendie, pas un thermomètre.

Les capteurs de la douleur ne se déclenchent que sur ce qui est fort: une chaleur au-delà de quarante degrés et quelques, une forte pression, des substances corrosives. Ils s'habituent beaucoup moins que les capteurs du toucher: une alarme ne doit pas se taire tant que la menace est là.

\section*{Deux fils}

Souvenez-vous de la fois où vous vous êtes cogné le doigt. D'abord, une douleur aiguë et précise: où et quand. Un instant plus tard, une douleur sourde, diffuse, durable: à quel point c'est grave. Ce sont deux fils différents. L'un est rapide, son signal arrive en quelques dizaines de millisecondes. L'autre est lent, son signal met de l'ordre d'une seconde.

\section*{La porte}

Pourquoi frotte-t-on l'endroit où l'on s'est cogné? Dans la moelle épinière, le signal de douleur passe par un \og portillon\fg{}. Le toucher réveille un neurone inhibiteur qui le referme en partie. Frotter la bosse aide donc vraiment, mais seulement contre une douleur modérée; une douleur forte passe quand même. Que la douleur forte coupe elle-même le neurone inhibiteur fait partie de la théorie initiale du portillon, mais n'a pas encore de confirmation directe.

\pencilfig{126}{%
% gate: the pain wire drives the relay neuron; touch wakes an inhibitory neuron that closes the gate
\draw[pen=pcAmber] (0,1.35) -- (3.05,1.0); \draw[arr=pcAmber] (2.8,1.03) -- (3.08,1.0);
\node[lbl, text=pcAmber!70!black, anchor=east] at (-0.05,1.35) {douleur};
\draw[pen=pcBlue] (0,-0.15) -- (1.75,-0.15); \draw[arr=pcBlue] (1.5,-0.15) -- (1.78,-0.15);
\node[lbl, text=pcBlue, anchor=east] at (-0.05,-0.15) {toucher};
\draw[dotpen=pcRed, wash=pcRed] (2.05,-0.15) circle (0.26);
\draw[pcRed, line width=0.7pt, -{Bar[width=6pt]}] (2.25,0.05) -- (3.2,0.66);
\node[lbl, text=pcRed, anchor=north] at (2.05,-0.45) {frein};
\draw[dotpen=pcGray, wash=pcGray] (3.4,0.9) circle (0.34);
\node[lbl, anchor=south] at (3.4,1.3) {portillon};
\draw[pen] (3.75,0.9) -- (5.6,0.9); \draw[arr] (5.35,0.9) -- (5.65,0.9);
\node[lbl, anchor=west] at (5.7,0.9) {vers le cerveau};
}{Le portillon dans la moelle épinière. Le signal de douleur gagne le cerveau par un neurone-portillon, et le toucher réveille un neurone inhibiteur qui le referme en partie. Frotter la bosse aide donc, mais seulement contre une douleur modérée.}

\section*{Un bouton de volume venu d'en haut}

Le cerveau sait lui-même rendre la douleur plus forte ou plus faible. D'en haut, depuis l'encéphale, des canaux radio descendent vers le portillon: \nt{SERO}, \nt{NORA} et des substances particulières, semblables aux analgésiques, que le cerveau sécrète lui-même. Dans une situation dangereuse, la douleur peut presque disparaître: qui fuit un danger n'a pas le temps de boiter. L'attente d'un soulagement atténue aussi la douleur, en partie par ces mêmes substances.

\section*{Quand l'alarme reste bloquée}

Une douleur répétée augmente la sensibilité: le même coup se met à faire plus mal. C'est aussi un apprentissage, une boucle qui s'auto-entretient pendant que la plaie guérit. Mais une boucle forte a un danger: elle peut \og rester bloquée\fg{} et continuer à sonner alors que la plaie est déjà guérie, comme le groupe de neurones du deuxième chapitre qui s'entretient tout seul.

\section*{Professeur et interruption}

La douleur est aussi un professeur: pour le \og mieux ou moins bien que prévu\fg{} de la dopamine, elle agit comme une récompense négative. Et elle interrompt l'activité en cours, comme une bouffée de noradrénaline: laisse tout tomber, occupe-toi du danger. Quant à la réaction la plus rapide, retirer la main d'un objet brûlant, elle se referme directement dans la moelle épinière grâce à cette même paire de muscles antagonistes du chapitre précédent, avant même que vous ayez eu le temps de sentir la douleur.

\fullversion{14}
